\section{問題設定と座標・評価対象}
\label{sec:problem_setting}

\subsection{入力と出力}
\label{sec:task}

本研究が解く検出課題は、単眼画像 \(I\in\R^{H\times W\times 3}\) から、
規制コートの物理的な点とカメラ幾何を同時に推定することである。
コートは規制寸法が既知であり、物理的な点は十四点
\(\{\mathbf{X}_k\}_{k=0}^{13}\)（以降 \KPName）として順序を固定する。
検出器の出力は次の四つである。

\begin{enumerate}[leftmargin=1.6em]
  \item \textbf{キーポイント}：各チャネルが一つの物理点へ対応する十四枚のヒートマップ。
  \item \textbf{領域分割}：コート面を背景と六種類の領域に分けるクラス写像。
  \item \textbf{線}：コート平面上の線幅付き領域を投影した二値マスクと、
  線を種類ごとに区別する意味ラベル。
  \item \textbf{カメラ}：カメラ中心 \(\mathbf{C}\)、
  カメラから正準コートへの回転 \(\Rot{\Canon}{\Cam}\)、焦点距離 \(f\)。
\end{enumerate}

キーポイントの順序は画像座標から推測しない。
どのチャネルがどの物理点かはデータ側の契約であり、
学習時には密な教師と姿勢教師が同じ順序を指していることを検証する。

\subsection{フレーム}
\label{sec:frames}

\(\T{b}{a}\) はフレーム \(a\) の座標をフレーム \(b\) へ写す同次変換とする。
\cref{tab:frames} に本研究が用いるフレームを示す。
再構成が出力する正規化シーンフレーム \(\Scene_n\) と
メートル系シーンフレーム \(\Scene_m\) の間には一様な尺度だけを置き、
尺度を剛体変換の中へ隠さない。
各コートは固有のメートル系フレーム \(\Court_i\) を持ち、
下流のデータ生成が受け取るのはシーンとコートの相互変換だけである。
変換の向き、姿勢の直列化順序、可視性の層は
付録~\ref{app:contracts} に完全な契約として記す。

\begin{table}[t]
  \centering
  \caption{\textbf{座標フレーム。}
  すべての剛体変換は正規直交・行列式 \(+1\) の同次行列として保存する。
  一様尺度はメートルアダプタが持ち、
  剛体変換の内部へは隠さない。
  }
  \label{tab:frames}
  \small
  \begin{tabularx}{\linewidth}{@{}p{0.20\linewidth}p{0.24\linewidth}X@{}}
    \toprule
    フレーム & 単位と軸 & 契約 \\
    \midrule
    正規化シーン \(\Scene_n\) & 正規化、右手系 &
      再構成が公開する座標。ガウス平均と復元カメラはここに住む。 \\
    メートルシーン \(\Scene_m\) & メートル、右手系 &
      \(\Scene_n\) の一様尺度変換。受理された全コートの共通の土台。 \\
    コート \(\Court_i\) & メートル。\(+X\) 右サイドライン、
      \(+Y\) 遠ベースライン、\(+Z\) 上 &
      規制寸法と物理点の番号付け。 \\
    カメラ \(\Cam\) & OpenCV。\(+x\) 右、\(+y\) 下、\(+z\) 前 &
      \(\T{\Scene_m}{\Cam}\) はカメラからシーン。
      投影はその逆変換を用いる。 \\
    正準コート \(\Canon\) & メートル、対象コート相対 &
      カメラ端の選択 \(S\in\{I,R_z(\pi)\}\) を適用した後のコートフレーム。
      姿勢教師の基準。 \\
    モデル画像 & 補正後画素 &
      キーポイントのロジット、主点、焦点距離、再投影残差を同じフレームで扱う。 \\
    \bottomrule
  \end{tabularx}
\end{table}


\subsection{評価対象}
\label{sec:evaluation_scope}

評価は検出と幾何に分ける。

\begin{description}[leftmargin=1.6em]
  \item[検出] 画像から得られる \KPName、領域、線の予測精度。
  実写の held-out validation と、合成の trajectory-group test を別の列として報告する。
  \item[幾何] 検出結果を規制コートへ登録したときの整合性。
  検出精度と登録精度を一つの数値へ潰さない。
\end{description}

実写側の source には独立した test split が存在しない。
本研究が使う実写データは train と validation のみであり、
学習はこの validation へ強く集中しやすい。
提案手法は、実会場の再構成から新しい視点を生成することで、
この集中を視点分布の設計へ置き換える。

\subsection{視点分布という中心変数}
\label{sec:view_variable}

本研究の中心変数は、学習に渡る画像の視点分布である。
ここで視点とは、カメラ中心、姿勢、
およびコートのどの部分が見えているか（可視範囲）の組を指す。

視点分布を動かすとき、次の三つを同時に動かさない。

\begin{enumerate}[leftmargin=1.6em]
  \item \textbf{画像数}：視点を増やせば画像も増える。
  同じ枚数のまま狭域・広域を比較する条件を別に用意する。
  \item \textbf{外観}：会場・時間帯・表面色を変えれば、視点以外も変わる。
  本研究の生成画像は一つの再構成シーンの外観を継承するため、
  外観の多様化は学習時の画像処理へ委ねる。
  \item \textbf{教師の種類}：ヘッドを増やせば教師も増える。
  視点分布の効果を主張する比較では教師構成を固定する。
\end{enumerate}

本稿が提示する実験計画は、この三つを順に固定しながら
視点分布だけを動かす条件を定義する（\cref{sec:experiments}）。
本稿の時点で測定済みの値は、
checkpoint 選択時に記録された合成 validation の値のみである。
